\section{Phân tích và gợi ý sửa đổi}

\subsection{Đánh giá tổng hợp}
Dựa trên các chỉ số và hình trực quan:
\begin{itemize}[leftmargin=1.5em]
    \item \textbf{Tính toàn vẹn dữ liệu rất tốt}: tỉ lệ ảnh hợp lệ gần như tuyệt đối.
    \item \textbf{Độ phủ nhãn đầy đủ}: 103/103 class trong metadata đều xuất hiện trong annotation.
    \item \textbf{Rủi ro chính}: mất cân bằng class theo kiểu long-tail và khác biệt phân bố kích thước ảnh giữa train và test.
\end{itemize}

Mức đánh giá phù hợp với kết luận trong notebook và tài liệu assess: dữ liệu tốt để huấn luyện, nhưng cần kỹ thuật giảm thiên lệch lớp phổ biến.

\subsection{Gợi ý sửa đổi cho pipeline huấn luyện}
\begin{enumerate}[leftmargin=1.5em]
    \item Dùng \textbf{class-weighted cross-entropy} hoặc \textbf{focal loss} để tăng tín hiệu cho lớp hiếm.
    \item Áp dụng \textbf{class-balanced sampler} hoặc oversample ảnh chứa lớp hiếm (đặc biệt nhóm dưới 10 ảnh).
    \item Theo dõi \textbf{per-class IoU} song song với mIoU tổng để tránh kết luận sai do lớp lớn áp đảo.
    \item Chuẩn hóa resize/crop đa tỉ lệ giữa train và test, giảm domain gap về kích thước ảnh.
    \item Đánh giá thêm theo nhóm lớp (nấm, rau, thịt...) để tìm hướng gộp lớp nếu cần rút gọn bài toán.
\end{enumerate}

\subsection{Gợi ý sửa đổi cho notebook phân tích}
\begin{enumerate}[leftmargin=1.5em]
    \item Chuẩn hóa notebook thành module có thể chạy một lệnh (ví dụ \texttt{scripts/generate\_figures.py}) để tái lập dễ hơn.
    \item Lưu toàn bộ biểu đồ tự động vào thư mục cố định thay vì chỉ hiển thị trong notebook.
    \item Tách cell xử lý dữ liệu, cell trực quan và cell tổng hợp kết luận để giảm trùng lặp code.
    \item Ghi log rõ ràng số file lỗi, danh sách class hiếm và cảnh báo domain gap trong phần cuối notebook.
    \item Xuất thêm file tổng hợp markdown hoặc PDF tự động để thuận tiện báo cáo định kỳ.
\end{enumerate}

\subsection{Ưu tiên triển khai đề xuất}
Thứ tự ưu tiên thực tế cho vòng huấn luyện tiếp theo:
\begin{enumerate}[leftmargin=1.5em]
    \item Kích hoạt class-weighted loss + theo dõi per-class IoU.
    \item Bổ sung sampler cân bằng lớp hiếm.
    \item Chuẩn hóa pipeline resize/crop train-test.
    \item So sánh lại baseline sau khi áp dụng các thay đổi trên.
\end{enumerate}
